\documentclass[11pt,a4paper]{article}
\usepackage[T1]{fontenc}
\usepackage{lmodern,amsmath,amssymb,amsthm,booktabs,microtype}
\usepackage[margin=24mm]{geometry}
\usepackage[colorlinks=true,urlcolor=blue,linkcolor=blue]{hyperref}
\newcommand{\Js}{\mathcal J_s}
\newcommand{\av}[1]{\langle #1\rangle}
\newcommand{\W}{\mathcal W}
\title{The missing linear component of the cubic singlet\\Exact charge-sector channels and their double-scaling weight}
\author{Alan Day \and Chengheng\thanks{AI research assistant (OpenAI).}}
\date{Local research note --- 6 November 2026}
\begin{document}\maketitle
\begin{abstract}
We calculate the first Gaussian coupling chaos of the raw Hermitian
cubic derivative singlet. It is not generally a conserved function
of the Hamiltonian. Its exact finite-size map is diagonal in
standard $U(N)$ irreducible tensor channels, with coefficients
given below by binomial expressions. After subtracting the
ensemble projection onto the Hamiltonian in each charge sector,
the residual has double-scaled norm
$64\mathcal J_s^6 e^{-8\lambda}$. This identifies a surviving
component omitted by the third-chaos projection and sets the
next spectral calculation.
\end{abstract}
\section{What is added to the cubic construction}
Use standard CAR, $H=H_s$, $L=[H,\cdot]$, and
\[
 A=\sum_i(Lc_i)^\dagger Lc_i,\quad
 T_3=\tfrac12\left(\sum_i c_i^\dagger L^3c_i+\mathrm{h.c.}\right).
\]
The exact identity $M_3=T_3+\tfrac32LA$ and the previously
constructed $q^3$ channel $R$ give
\begin{equation}
 \boxed{M_3=R+aS+bZ+\bar t(Q_f)H+F+\tfrac32LA,}       \label{decomp}
\end{equation}
where $S=\W_3\{H,A\}/2$, $Z=\W_3H^3$,
$R=\W_3T_3-aS-bZ$, and
\[
 a=-\frac{(1-q)^2}{1+q^2},\qquad
 b=\frac{2\lambda(1-q)(1+q^3)}{(1+q^2)(1+q+q^2)},
 \quad q=e^{-4\lambda}.
\]
The new terms are defined exactly at finite $N$:
\begin{equation}
 \W_1T_3=\bar t(Q_f)H+F,\qquad
 \mathbb E_J\operatorname{tr}_m(HF)=0.                \label{F}
\end{equation}
Here $\operatorname{tr}_m$ is the normalized trace on charge $m$;
$\bar t(m)$ is a deterministic projection coefficient.
It follows that $\bar t(Q_f)H$ is conserved in each realization.
$F$ can still contain static components; (\ref{F}) is not a
projection onto the entire commutant.

In the full $\beta=\mu_{\rm chem}=0$ annealed trace, first take
$N,p\to\infty$ at $\lambda=p^2/N>0$, with
$\av{H^2}\to\sigma^2=\Js^2/\lambda$. Our norm result is
\begin{equation}
 \boxed{\lim_{\rm DS}\av{F^2}=64\Js^6q^2>0.}         \label{norm}
\end{equation}
The earlier $R$ has $\av{R^2}\sim32\Js^6/\lambda$ as
$\lambda\to0^+$, so $\av{F^2}/\av{R^2}\sim2\lambda$.
The raw derivative family therefore has an additional component
at the first relative order in this limit.

\newpage
\section{Exact coefficients from covariance channels}
Write the coupling covariance in standard CAR units as
$\mathbb E J_{IJ}J_{KL}=w\delta_{IL}\delta_{JK}$.
The normal-ordered $p$-body Hamiltonian decomposes into the
established, multiplicity-free $U(N)$ tensor sectors
\[
 H_m=\sum_{\nu=0}^{\min(p,m,N-m)}H_m^{(\nu)} .
\]
Define the following quantities; every binomial is zero outside
its allowed range:
\begin{align}
 C_m&=w\binom mp\binom{N-m+p}{p},\\
 \kappa_\nu(m)&=w\binom{m-\nu}{p}\binom{N-m+p-\nu}{p},\\
 j_c(m)&=w\binom{m-1}{p}\binom{N-m+p}{p},\\
 j_-(m,\nu)&=w\binom{m-\nu}{p}\binom{N-m+p-\nu+1}{p},\\
 j_+(m,\nu)&=w\binom{m-\nu-1}{p}\binom{N-m+p-\nu}{p},\\
 \ell_\nu&=\frac{\nu(N-\nu-p+1)}{N-2\nu+1}.
\end{align}
$C_m=\mathbb E K_m^2$ is scalar, and
$\mathbb E K_m H_m^{(\nu)}K_m=\kappa_\nu(m)H_m^{(\nu)}$
for an independent copy $K$ of $H$. These covariance eigenvalues
are standard embedded-ensemble results \cite{TT,Kota}.
The $j_\pm$ are the corresponding particle-removal channels.

The first Gaussian chaos is
\begin{equation}
 (\W_1T_3)_m=\sum_\nu t_\nu(m)H_m^{(\nu)},             \label{map}
\end{equation}
with the exact coefficient
\begin{align}
 t_\nu(m)={}&-(m+2p)C_{m-1}+(m-3p)C_m+3p\,j_c(m)
 \nonumber\\
 &+(m-p)\kappa_\nu(m-1)-m\kappa_\nu(m)
 \nonumber\\
 &+3\{\ell_\nu j_-(m,\nu)+(p-\ell_\nu)j_+(m,\nu)\}.
                                                               \label{t}
\end{align}
This result keeps the finite-$N$ rank dependence. Replacing all
$\kappa_\nu/C_m$ by the common leading crossing factor too early
would erase the residual (\ref{norm}).

For example, at $(N,p,m)=(5,2,3)$ and $w=1$,
\[
 (t_0,t_1,t_2)=(-96,-167/2,-88).
\]
Distinct coefficients on noncentral sectors allow nonconservation.
A deterministic Gaussian-integer Hamiltonian specified with the
saved validation data gives the exact witness
\begin{equation}
 \operatorname{tr}_3
 [H,\W_1T_3]^\dagger[H,\W_1T_3]=\frac{8933634}{5}>0. \label{witness}
\end{equation}
Since the subtraction in (\ref{F}) commutes with $H$, the same
witness applies to $[H,F]$. It tests a finite system, not a
double-scaling dynamical limit.

\newpage
\section{Derivation through particle-removal tensors}
For a cubic polynomial $P$, $\W_1P(H)$ is the sum of its three
terms containing two copies of an independent Gaussian $K$ and
one $H$, averaged over $K$. Apply this to
\[
 L^3c=H^3c-3H^2cH+3HcH^2-cH^3 .
\]
Let
$\mathcal J_m(Y)=\mathbb E K_{m-1}YK_m$ and
$\Phi(Y)=\sum_i c_i^\dagger Yc_i$.
Homogeneity gives $\Phi(H_{m-1})=(m-p)H_m$ and
$\mathcal J_m(c_i)=j_c(m)c_i$. Grouping the three contractions
in each monomial gives
\begin{align}
 (\W_1M_3)_m={}&
 [-(m+2p)C_{m-1}+(m-3p)C_m+3p j_c]H_m \nonumber\\
 &+\Phi(\mathbb E K_{m-1}H_{m-1}K_{m-1})
      -m\mathbb E K_mH_mK_m \nonumber\\
 &+3\sum_i c_i^\dagger\mathcal J_m([c_i,H]).
                                                               \label{operator}
\end{align}
For a traceless charged tensor with $r$ creations and $s$ annihilations,
input charge $m$, the covariance multiplier is
\begin{equation}
 \mathcal J(Y)=
 w\binom{m-s}{p}\binom{N-m+p-r}{p}\,Y.                \label{charged}
\end{equation}
It follows by normal ordering a highest-weight disjoint-index
tensor: the paired $K$ indices have $m-s$ available occupied
slots and $N-m+p-r$ available complementary slots. Equivariance
extends the result to its irreducible tensor space.

To resolve $[c_i,H^{(\nu)}]$, write
$H_m^{(\nu)}=f_m O_\nu$ with
$f_m=\binom{m-\nu}{p-\nu}$ and $O_\nu$ a primitive traceless
$\nu$-body tensor. Set $u_i=c_iO_\nu$, $v_i=O_\nu c_i$.
The difference $u_i-v_i$ belongs to the lower charged sector
$(r,s)=(\nu-1,\nu)$. Only this and the upper sector
$(\nu,\nu+1)$ occur; multiplicity freedom makes the lower
projection proportional to $u-v$.
In unnormalized sector trace, dividing
all inner products by $\operatorname{Tr}_m O_\nu^\dagger O_\nu$,
\[
 \|u\|^2=m,\quad (u,v)=m-\nu,\quad
 \|v\|^2=\frac{(N-m+1)(m-\nu)}{N-m-\nu+1}.
\]
Thus the lower projection of $u$ is
$(N-m-\nu+1)(u-v)/(N-2\nu+1)$.
Since $f_{m-1}/f_m=(m-p)/(m-\nu)$, contraction with
$\sum_i c_i^\dagger$ gives weights
$\ell_\nu$ and $p-\ell_\nu$ in the lower and upper charged
sectors. Using (\ref{charged}) in (\ref{operator}) proves
(\ref{t}). For $\nu=0$ the lower weight is zero, understood
without a negative-rank tensor.

All coefficients are real. The resulting map preserves Hermiticity,
so $\W_1M_3=\W_1T_3$ and $\W_1LA=0$. The displayed derivation
therefore also checks that the exact time-descendant subtraction
does not hide another first-chaos term.

\newpage
\section{Rank weights and the surviving norm}
Let $d_\nu=\binom N\nu^2-\binom N{\nu-1}^2$, and use the
standard covariance eigenvalue notation \cite{Kota}
\[
 \Lambda^\nu(N,m,r)=\binom{m-\nu}{r}
                         \binom{N-m+r-\nu}{r}.
\]
The normalized Hamiltonian weight of channel $\nu$ is
\begin{equation}
 \omega_\nu=
 \frac{d_\nu\Lambda^\nu(N,m,m-p)}
      {\binom Nm\Lambda^0(N,m,p)},\qquad \sum_\nu\omega_\nu=1 .
\end{equation}
Different tensor components are independent Gaussian variables
and orthogonal in the ensemble Hilbert--Schmidt inner product.
Consequently
\begin{align}
 \bar t(m)&=\sum_\nu\omega_\nu t_\nu(m),\\
 \mathbb E\operatorname{tr}_m F^2
 &=C_m^3\,\operatorname{Var}_{\omega}
          \left(\frac{t_\nu(m)}{C_m}\right).           \label{variance}
\end{align}

At half filling $m=N/2$, put $s=p-\nu$. The exact adjacent-weight
ratio is
\begin{equation}
 \frac{\omega_{\nu-1}}{\omega_\nu}
 =\left(\frac{\nu}{N-\nu+1}\right)^2
 \frac{1-\left(\frac{\nu-1}{N-\nu+2}\right)^2}
      {1-\left(\frac{\nu}{N-\nu+1}\right)^2}
 \frac{m-\nu+1}{p-\nu+1}
 \frac{N-p-\nu+1}{N-m-\nu+1}.                         \label{ratio}
\end{equation}
For fixed $s$ this tends to $\lambda/(s+1)$. For sufficiently
large $N$ and all $0\le s<p$, it is bounded by
$32\lambda/(s+1)$. Normalization and this factorial tail give
\begin{equation}
 s\ \Longrightarrow\ {\rm Poisson}(\lambda)
\end{equation}
with convergence of every polynomial moment.

Writing $c_\nu=t_\nu/C_m$, direct expansion of (\ref{t}) gives
\begin{equation}
 c_{p-s}-c_p\longrightarrow 8\lambda q\,s.             \label{slope}
\end{equation}
One way to see the coefficient is to write $c_\nu=B_0+
(\kappa_\nu/C_m)E_\nu$, where $B_0$ is independent of $\nu$.
At half filling
$E_p/p\to2$, $E_{p-s}-E_p=O_s(p/N)$,
$\kappa_p/C_m\to q$, and
$\kappa_{p-s}/\kappa_p=1+4\lambda s/p+O_s(p^{-2})$.
The rational factors in (\ref{t}) also bound
$|c_{p-s}-c_p|\le C_\lambda(1+s)$ uniformly for $s\le p$.
Together with the factorial tail, this justifies taking the
variance, rather than only a pointwise coefficient limit:
\begin{equation}
 \operatorname{Var}_\omega(c_\nu)\longrightarrow
                  64\lambda^3q^2 .                  \label{varlimit}
\end{equation}

\newpage
\section{Full trace, scale comparison, and numerical checks}
The half-filled result extends to the full
$2^{-N}\operatorname{Tr}$ state. For $m/N$ in a compact
neighborhood of $1/2$, the ratio and slope estimates above are
uniform, with continuous density dependence.
The binomial charge distribution concentrates at $1/2$.
Also $C_m$ is maximal near $m=(N+p)/2$ and
$C_{\max}/C_{N/2}=\exp[O(p^3/N^2)]\to1$.
Hence $C_m/\sigma^2\to1$ on the concentrating charge window.
Outside a fixed density window the binomial charge tail is
exponentially small, while (\ref{t}) bounds the coefficients
by $O(N C_{\max})$; these tails do not contribute to
(\ref{variance}). Equations (\ref{variance}) and (\ref{varlimit})
then give (\ref{norm}).

The cubic-chaos result from the preceding stages is
\[
 \av{R^2}_{\rm DS}=\sigma^6V,\qquad
 V=\frac{4\lambda(1+q)[1-q^3+4\lambda q^2(2q^2-q+2)]}
 {(1+q^2)(1+q+q^2)}.
\]
Therefore the exact double-scaled equal-time ratio is
\begin{equation}
 \frac{\av{F^2}_{\rm DS}}{\av{R^2}_{\rm DS}}
 =\frac{64\lambda^3q^2}{V}=2\lambda+O(\lambda^2).
\end{equation}
Gaussian chaos makes the linear and cubic components orthogonal
at equal time. Real-time evolution depends on the same random
$H$ and does not preserve that chaos decomposition; their
time-dependent cross correlators must be computed.

The channel formula was checked against direct Gaussian
polarization in occupation-basis CAR matrices for
$(N,p)=(3,1),(4,2),(5,2),(6,2),(6,3)$: 42 allowed sector/rank
coefficients agree exactly with bounded Gaussian-integer
arithmetic. Independent tests cover 31 charged covariance
eigenvalues, 199 exact rational weight normalizations and ratios,
three symbolic projection identities, and the nonzero witness
(\ref{witness}).

The full rank sum, with no rank cutoff, was evaluated at 80 and
110 decimal digits for three $\lambda$ values and
$p=16,32,64,128$. The table shows the relative difference
of the variance from (\ref{varlimit}):
\begin{center}
\begin{tabular}{rrrrr}
\toprule
$\lambda$ & $p=16$ & $p=32$ & $p=64$ & $p=128$\\
\midrule
0.10 & $-13.97\%$ & $-6.94\%$ & $-3.45\%$ & $-1.72\%$\\
0.25 & $-7.76\%$ & $-3.51\%$ & $-1.66\%$ & $-0.81\%$\\
0.50 & $+3.54\%$ & $+2.50\%$ & $+1.41\%$ & $+0.74\%$\\
\bottomrule
\end{tabular}
\end{center}
The finite-$p$ differences are retained rather than treated as
numerical error. Precision comparisons are below $10^{-55}$ in
relative variance. The limiting statement is supported by the
ratio and uniform-bound argument, not inferred from this table.

\newpage
\section{Physical interpretation and the next calculation}
The original derivative operator is now decomposed into explicit
pieces in (\ref{decomp}). The known $q^3$ channel remains intact.
The missing linear part has an exactly known finite-$N$ map;
its energy-orthogonal remainder survives double scaling in norm,
and finite-size examples demonstrate nonconservation.

The next question is its \emph{dynamical} spectral measure.
A positive equal-time residual alone does not establish a new
propagating particle or rule out additional conserved projections.
In particular, identifying a $q$-probe kernel for $F$ requires
same-coupling mixed moments or an explicit covariance reduction.
It cannot be concluded merely because $\nu/p\to1$.
The rank-dependent terms in (\ref{t}) show exactly which
subleading coefficients must be kept in that calculation.

The concrete next step is to evaluate
$\mathbb E\operatorname{tr}(F L^2F)$ and the two-energy strength
density using the resolved covariance channels, including
$H$--$F$ cross contractions. This will determine how the residual
fits into the full raw-$M_3$ correlator. We can then compute mixed
correlators with the already established $d=2$ and $d=3$ singlets,
to supply new interaction data for the Susskind program.
Their physical mass assignment still needs the spatial dictionary.

\paragraph{Source boundary.}
The irreducible decomposition, covariance eigenvalues, dimensions
and channel weights used here come from embedded-ensemble
representation theory \cite{Kota,TT}. The charged multiplier
was obtained by the same normal-order counting and explicitly
checked. The new calculation recorded in this note is the
application to $\W_1T_3$, its finite-sector coefficients,
and the energy-orthogonal residual limit.
The source papers' spectral-density correlation functions are
not confused with the transition-strength correlator of $F$.
No priority claim is made for standard representation theory.

All scripts and raw outputs are local. Earlier cubic-kernel
and other scientific artifacts were preserved byte-for-byte.
The six-page PDF was actually compiled and visually checked.
No Git commit, push or archive was made; the complete
de Sitter research objective remains active.

\begin{thebibliography}{9}
\bibitem{Kota}
V. K. B. Kota,
\emph{Bivariate moments of the two-point correlation function
for embedded Gaussian unitary ensemble with $k$-body interactions},
\href{https://arxiv.org/html/2208.11312}{arXiv:2208.11312v2},
Sec.~IV, Eqs.~(28)--(39).
\bibitem{TT}
J. Tall and S. Tomsovic,
\emph{Double-scaled bosonic and fermionic embedded ensembles,
complex SYK, and the dual Hilbert space},
\href{https://arxiv.org/html/2604.14522v2}{arXiv:2604.14522v2},
Sec.~2.1, Eqs.~(2.8)--(2.19).
\end{thebibliography}
\end{document}
